%%% Не мянять - Do not modify
%%
%%
\clearpage                                  % В том числе гарантирует, что список литературы в оглавлении будет с правильным номером страницы
%\hypersetup{ urlcolor=black }               % Ссылки делаем чёрными
%\providecommand*{\BibDash}{}                % В стилях ugost2008 отключаем использование тире как разделителя 
\urlstyle{rm}                               % ссылки URL обычным шрифтом
\ifdefmacro{\microtypesetup}{\microtypesetup{protrusion=false}}{} % не рекомендуется применять пакет микротипографики к автоматически генерируемому списку литературы
%\newcommand{\fullbibtitle}{Список литературы} % (ГОСТ Р 7.0.11-2011, 4)
%\insertbibliofull  
%\noindent
%\begin{group}
\chapter*{Список использованных источников}	
\begin{enumerate}[1. ]
    \item Paygude P., Joshi S.~D., Joshi M. Fault Aware Test Case 
    Prioritization in Regression Testing using Genetic Algorithm~// 
    International Journal of Emerging Trends in Engineering Research. 
    --- 2020. --- Vol.~8, No.~5. --- P.~2112--2117. --- 
    DOI:~10.30534/ijeter/2020/104852020.

    \item Mahdieh M., Mirian-Hosseinabadi S.-H., Etemadi K., 
    Bohlouli M. Incorporating Fault-Proneness Estimations into 
    Coverage-Based Test Case Prioritization~// 
    Information and Software Technology. --- 2021. --- Vol.~133. --- 
    P.~106483. --- DOI:~10.1016/j.infsof.2021.106483.

    \item Elbaum S., Malishevsky A.~G., Rothermel G. Test Case 
    Prioritization: A Family of Empirical Studies~// 
    IEEE Transactions on Software Engineering. --- 2002. --- 
    Vol.~28, No.~2. --- P.~159--182. --- 
    DOI:~10.1109/32.988497.

    \item Руденко~М.~А. tcp-methods: репозиторий данных для 
    приоритизации тестовых случаев~// GitHub. --- URL: 
    https://github.com/oneren38/tcp-methods (дата обращения: 
    08.04.2026).

    \item dathpo. Test Case Prioritisation --- Genetic Algorithm: 
    BigFaultMatrix~// GitHub. --- URL: 
    https://github.com/dathpo/Test\_Case\_Prioritisation\_-\_Genetic\_Algorithm 
    (дата обращения: 08.04.2026).
\end{enumerate}

\label{references}
\addcontentsline{toc}{chapter}{Список использованных источников}	% в оглавление 
% \printbibliography[env=SSTfirst]                         % Подключаем Bib-базы
%\ifdefmacro{\microtypesetup}{\microtypesetup{protrusion=true}}{}
%\urlstyle{tt}                               % возвращаем установки шрифта ссылок URL
%\hypersetup{ urlcolor={urlcolor} }          % Восстанавливаем цвет ссылок



%\urlstyle{rm}                               % ссылки URL обычным шрифтом
%\ifdefmacro{\microtypesetup}{\microtypesetup{protrusion=false}}{} % не рекомендуется применять пакет микротипографики к автоматически генерируемому списку литературы
%\insertbibliofull                           % Подключаем Bib-базы
%\ifdefmacro{\microtypesetup}{\microtypesetup{protrusion=true}}{}
%\urlstyle{tt}                               % возвращаем установки шрифта ссылок URL
